This family looks inside invalid plans. Instead of asking only whether VAL accepted the final sequence, it asks how far the plan executes and what kind of state-tracking error causes the first failure.

\begin{table}[H]\small
\centering
\begin{tabularx}{\linewidth}{@{}l l Y Y@{}}
\toprule
\textbf{Metric} & \textbf{Formula} & \textbf{What it measures} & \textbf{Main interpretation}\\
\midrule
Exec & \(\mathrm{first\_failure}/\mathrm{plan\_length}\) & Prefix fraction that executes before the first failure. & High value means the plan respects local transitions for many steps.\\
\addlinespace
Sequencing errors & Failed preconditions that were true earlier. & Whether the model had the needed state but misplaced it. & Ordering failure rather than total fabrication.\\
\addlinespace
State fabrications & Failed preconditions that were never true. & Whether the model required an impossible state. & Weak or incoherent world modelling.\\
\addlinespace
PAS & \(\frac{\mathrm{sequencing}}{\mathrm{sequencing}+\mathrm{fabrication}}\) & Share of failures that are ordering mistakes. & High PAS means failures are closer to valid causal reasoning.\\
\addlinespace
Temporal distance & \(\mathrm{failure\_step}-\mathrm{last\_true\_step}\) & Distance from last valid prerequisite to failure. & Large distance means long-horizon ordering collapse.\\
\bottomrule
\end{tabularx}
\caption{Execution and reasoning metrics derived from simulated replay.}
\label{tab:execution-metrics-compact}
\end{table}

These metrics should be read as a cascade. Executability is the first gate; PAS becomes meaningful only when the plan runs far enough for the simulator to observe state transitions. For near-solving models, very small failure counts make sequencing/fabrication proportions unstable, so the breakdown must be read with SR and FASR.